% SPDX-License-Identifier: Apache-2.0
% covers: Sd
\datasheet{Sd}{A native-mode SD card host on AXI-Lite: one command and
at most one block at a time, the block through a buffer of 128 words.}

\dsfacts{\code{lib/parts/src/sd.rs}}{\code{txhdl\_parts::sd::Sd}}%
{\code{//docs:examples}}

\subsection*{Function}

Operating an SD card in SPI mode limits throughput to several hundred
kilobytes per second over a single data line. In contrast, SD native
mode uses a bidirectional protocol over dedicated physical lines: a command
line, \code{CMD}, conveys 48-bit commands and returns 48-bit or 136-bit
responses. Data transfers use one or four data lines (\code{DAT}) to
transfer 512-byte blocks, protected by per-line CRC-16 checksums. Outgoing
commands include a CRC-7 checksum, as do all responses with the exception
of \code{ACMD41}, whose CRC field is defined as all ones.

\code{Sd} implements this native host controller behind an AXI-Lite port,
\code{bus}. Software configures the command argument and writes the command
control word. The command word selects the transfer parameters: response
type (none, short, or long), data transfer direction (read block or write
block), busy line monitoring, and CRC checking. The controller dispatches
the command, samples the response, transfers block data through its internal
buffer via the \code{data} register, and updates execution state in \code{status}.

Multi-block transfers issue an initial command (\code{CMD18} or \code{CMD25}),
stream data blocks across subsequent cycles, and terminate with \code{CMD12}.
Because streaming reads at 25~MHz overloads register polling loops, the host
integrates dedicated DMA paths to system memory (issue 912). Two DMA engines
from issue 151 interface with the controller: a \code{LineStore} instance
for reads and a \code{LineFetch} instance for writes. Setting a non-zero count
in \code{blocks} initiates automated multi-block DMA: the internal buffer
acts as a ring buffer between physical SD lines and the DMA engine. Read data
streams to the store engine as the card presents bytes; write data streams
from the fetch engine, dispatching each block once fully buffered. The transfer
completes only after the DMA engine retires. Commands perform either reads or
writes, never both simultaneously, ensuring mutual exclusion between the DMA
engines.

The card's side is eight wires: \code{sclk}, the command line in, out
and its output enable, the four data lines in, out and their output
enable, and \code{irq}. The engines' side is nine more:
\code{dma\_in} and \code{dma\_out}, the words from the fetch engine
and to the store engine; \code{dma\_at}, \code{dma\_bytes} and
\code{dma\_words}, where in memory and how much; \code{store\_go}
and \code{fetch\_go}, the starts; and \code{store\_busy} and
\code{fetch\_busy}, the engines' busy lines.

\subsection*{Parameters}

None.

\subsection*{Register map}

The word is address bits 5 to 2. A word the map does not name reads as
zero.

\dsregs{Sd}

\code{cmd}'s \code{resp} is 0 for none, 1 for short and 2 for long.
Writing \code{cmd} starts the command. Writing a one to
\code{status}'s \code{done} clears it and the four faults with it.
\code{dma} is where in memory the next command's blocks start, and
\code{blocks}' \code{count} how many it moves, zero for the buffer as
before; \code{left} counts the blocks still to move and \code{run}
says an engine is still moving words.

A half of a card clock takes \code{div + 1} cycles, so the card's
clock is the system's over \code{2 * (div + 1)}: at 100~MHz, 124 is
the 400~kHz a card starts at, 1 is 25~MHz, the fastest this host
runs, and 0 runs as 1 (issue 929).

\dstables{Sd}

\subsection*{Behaviour}

\begin{itemize}
\item \textbf{A write to \code{cmd} while a command runs is ignored.}
A program reads \code{status} first.
\item The host drives its lines on the falling edge of the clock it
makes and samples the card's on the rising edge. The clock runs only
while a command runs, and is low otherwise.
\item A response that has not started 64 card clocks after the command
ends is a timeout. A block, the CRC status token after a written
block, and the end of the card's busy are each given $2^{24} - 1$ card
clocks.
\item A CRC is checked the way a shift register checks one: the
received CRC goes into the same register after the data, and what is
left is zero when the two agree. For a long response the CRC covers
the 120 bits after the header. With four lines each line has its own
CRC-16, and any of the four being wrong sets bit 5.
\item A written block is accepted only when the card's CRC status
token is \code{010}; anything else sets bit 5. A command whose
response CRC is wrong goes no further: its block does not move.
\item Starting a read sets the write pointer to zero, and starting a
write sets the read pointer to zero, so a block fills or empties the
buffer from its start. A read of \code{data} moves the read pointer
and a write moves the write pointer, each modulo 128.
\item \code{irq} is done and the interrupt enable together, so it
stays high until the program clears done.
\item Every AXI-Lite access is answered \code{OKAY}.
\end{itemize}

\subsection*{Verification}

\begin{itemize}
\item \code{ex\_sd} does what a driver does, through an AXI4 host and
\code{LiteBridge<1, ..>}, against \code{SdCard}, the model of a card in \code{sd.rs},
stepped on the wires every cycle. It brings the card up with
\code{CMD0}, \code{CMD8}, \code{ACMD41}, \code{CMD2}, \code{CMD3},
\code{CMD7} and \code{ACMD6}, reads block 3 on one line, then writes
block 5 on four and reads it back, and checks both blocks.
\item The netlist is checked against that run under nvc and
Verilator: \code{//docs:sd\_sim\_sd\_tb\_test} and
\code{//docs:sd\_vsim\_test}.
\item \code{ex\_sddma} moves blocks through memory with real
\code{LineStore<32, 1, 16, 16>} and \code{LineFetch<32, 1, 16, 16>}
engines behind an \code{Arbiter2}, the card sending its blocks back to
back, and its netlist is checked against that run:
\code{//docs:sddma\_sim\_sd\_dma\_tb\_test} and
\code{//docs:sddma\_vsim\_test}.
\item \code{//lib/parts:parts\_test} holds eleven tests of it:
\begin{itemize}
\item the card comes up and gives its CID;
\item a block is read on one line;
\item a block is written and read back on four lines, the card's CRC
status accepting it;
\item a card slow to drive is read and written at 25~MHz;
\item a divider of zero runs as one;
\item several blocks each way, as a command, data-only transfers and
\code{CMD12};
\item a spoiled response CRC sets bit 3 and the next command is
clean; a command the card does not answer sets bit 2; a spoiled block
sets bit 5;
\item a card wants its clocks before its first command;
\item the CRC-7 of \code{CMD0} with a zero argument is \code{0x4a} and
the CRC-16 of 512 bytes of \code{0xff} is \code{0x7fa1}, the
specification's own examples;
\item blocks read back to back go to memory through the store engine;
\item blocks written from memory through the fetch engine reach the
card.
\end{itemize}
\end{itemize}

\subsection*{Used by}

The example \code{ex\_sd}, and the board: the Vreteno board holds it
in the peripheral page's ninth slot, at \code{0x3800}, with the card
slot's pins in bank 16 (E13 the clock, E14 the command line, D15,
D14, F14 and F13 the data lines) on every top that holds the board.
Its engines are the board's \code{sdstore} and \code{sdfetch}, sharing
one host, \code{sdshost}, through \code{sdarb}, an
\code{Arbiter2<32, 32, 4, 1, 2, 0>}, on the main arbiter's sixth port.
\code{cpu/vreteno/rust/sdprobe.rs} brings a card up and prints its CID
and block 0, then reads blocks 0 to 3 with one \code{CMD18} straight
into the DDR3 and compares them with the same blocks read one at a time
through the buffer; \code{//cpu/vreteno:board\_test} runs it against
\code{SdCard} and expects \code{dma same}, and
\code{docs/board-checks.md} says how to run it on the board.

\subsection*{Limits}

Only reading has met a real card: on 2026-10-03 the board read a
32~GB card's CID and block 0, its partition table, twice (issue 153),
and on 2026-10-04 read them again at 25~MHz, line for line the same
(issue 929).
No write has been sent to a card on the board, and \code{sdprobe}
sends none.

One command at a time. Without a count in \code{blocks} a command
moves at most one block, through a register; with one, its blocks go
through memory by the engines, which only a read or a write uses at a
time.

\code{SdCard}, the model, is not a component: nothing lowers it, and
it is only as right as the tests that check the host against it.

It runs on one clock, \code{DefaultClock}.
